
\documentclass[addpoints]{exam}
\usepackage{amsmath}
\usepackage{amssymb}
\usepackage{color}
\usepackage{siunitx}

% \printanswers

\newcommand{\event}{Trigonometric Functions}
\newcommand{\tournament}{}
\def\type{0}

\setlength\fillinlinelength{\textwidth}
\CorrectChoiceEmphasis{\color{red}\bfseries}
\SolutionEmphasis{\color{red}\bfseries}

\setlength\answerclearance{2pt}
\pagestyle{headandfoot}
\runningheadrule
\runningheader{\event}{\tournament}{\ifnum\type=1 Class Set Test \else Full Name:\hspace{1cm} \fi}
\runningfootrule
\runningfooter{}{Page \thepage\ of \numpages}{}

\begin{document}
\begin{center}
    {\huge Grade 11 Functions \\ \event
    \ifprintanswers \space Key
    \else \space \ifnum\type<2 Test \fi \ifnum\type=1 \textbf{(CLASS SET)} \fi
          \ifnum\type=2 Answer Sheet \fi
    \fi \\ \vspace{3mm}\tournament\vspace{1mm}}
\end{center}

\vspace{.25\textheight}

\begin{center}
    First Name: \ifprintanswers
    \underline{\hspace{3.5cm}}\textbf{KEY}\underline{\hspace{5.5cm}}\\\vspace{5mm}
    \else \underline{\hspace{10cm}}\\ \vspace{5mm} \fi
    Last Name: \underline{\hspace{10cm}}\\ \vspace{5mm}
\end{center}

\noindent\textbf{\underline{Directions:}}
\begin{itemize}
    \ifnum\type=1 \item \textbf{This test is a class set.} Do not write on this paper. \fi
    \item Show all work for full marks. Angles in radians unless stated.
    \item The test is designed to be completed in 75 minutes.
\end{itemize}
\begin{center}
\textit{For grading use only}\\
\multirowgradetable{1}[pages]
\end{center}

\newpage
\begin{questions}

\section*{Part A --- Multiple Choice (15 marks)}
\vspace{5mm}

\question[1] Convert $\ang{270}$ to radians.
\begin{choices}
    \choice $\pi$
    \choice $2\pi$
    \correctchoice $\dfrac{3\pi}{2}$
    \choice $\dfrac{\pi}{2}$
\end{choices}

\question[1] Convert $\dfrac{5\pi}{6}$ to degrees.
\begin{choices}
    \choice $\ang{30}$
    \choice $\ang{120}$
    \correctchoice $\ang{150}$
    \choice $\ang{300}$
\end{choices}

\question[1] The period of $y = \sin x$ is:
\begin{choices}
    \choice $\pi$
    \correctchoice $2\pi$
    \choice $\dfrac{\pi}{2}$
    \choice $4\pi$
\end{choices}

\question[1] The amplitude of $y = -3\sin x$ is:
\begin{choices}
    \choice $-3$
    \correctchoice $3$
    \choice $\dfrac{1}{3}$
    \choice $6$
\end{choices}

\question[1] The period of $y = \cos(2x)$ is:
\begin{choices}
    \choice $4\pi$
    \correctchoice $\pi$
    \choice $2\pi$
    \choice $\dfrac{\pi}{2}$
\end{choices}

\question[1] The graph of $y = \sin\!\left(x - \dfrac{\pi}{4}\right)$ is shifted from
$y = \sin x$ by:
\begin{choices}
    \choice $\dfrac{\pi}{4}$ to the left
    \correctchoice $\dfrac{\pi}{4}$ to the right
    \choice $\dfrac{\pi}{4}$ upward
    \choice $\dfrac{\pi}{4}$ downward
\end{choices}

\question[1] The equation $y = 2\sin(x) + 1$ has:
\begin{choices}
    \choice Amplitude 1, vertical shift 2
    \correctchoice Amplitude 2, vertical shift up 1
    \choice Amplitude 2, vertical shift down 1
    \choice Period $\pi$, vertical shift 1
\end{choices}

\question[1] What is the range of $y = 4\cos x - 1$?
\begin{choices}
    \choice $[-4, 4]$
    \choice $[-1, 1]$
    \correctchoice $[-5, 3]$
    \choice $[0, 4]$
\end{choices}

\question[1] The maximum value of $y = -2\sin(3x) + 5$ is:
\begin{choices}
    \choice $5$
    \choice $-2$
    \correctchoice $7$
    \choice $3$
\end{choices}

\question[1] The period of $y = \sin\!\left(\dfrac{x}{3}\right)$ is:
\begin{choices}
    \choice $\dfrac{2\pi}{3}$
    \choice $\pi$
    \correctchoice $6\pi$
    \choice $2\pi$
\end{choices}

\question[1] A sinusoidal function has amplitude 4 and period $\pi$. Which equation
could represent it?
\begin{choices}
    \choice $y = 4\sin x$
    \correctchoice $y = 4\sin(2x)$
    \choice $y = 4\sin\!\left(\dfrac{x}{2}\right)$
    \choice $y = 2\sin(4x)$
\end{choices}

\question[1] The arc length of a sector with radius 6 and central angle $\dfrac{2\pi}{3}$
rad is:
\begin{choices}
    \choice $\dfrac{2\pi}{3}$
    \choice $2\pi$
    \correctchoice $4\pi$
    \choice $8\pi$
\end{choices}

\question[1] $\cos\!\left(\dfrac{3\pi}{2}\right) =$
\begin{choices}
    \choice $1$
    \choice $-1$
    \correctchoice $0$
    \choice $\dfrac{\sqrt{2}}{2}$
\end{choices}

\question[1] The graph of $y = \tan x$ has vertical asymptotes at:
\begin{choices}
    \choice $x = n\pi$, $n \in \mathbb{Z}$
    \correctchoice $x = \dfrac{\pi}{2} + n\pi$, $n \in \mathbb{Z}$
    \choice $x = 2n\pi$, $n \in \mathbb{Z}$
    \choice $x = \dfrac{\pi}{4} + n\pi$, $n \in \mathbb{Z}$
\end{choices}

\question[1] An arc of length $9\pi$ cm on a circle of radius 6 cm subtends an angle of:
\begin{choices}
    \choice $\dfrac{3\pi}{2}$ rad
    \correctchoice $\dfrac{3\pi}{2}$ rad
    \choice $\pi$ rad
    \choice $2\pi$ rad
\end{choices}


\section*{Part B --- Short Answer (33 marks)}

\question Radian measure and the unit circle.
\begin{parts}
    \part[2] Convert $\ang{225}$ to radians and $\dfrac{7\pi}{4}$ to degrees.
    \part[2] Find the exact values of $\sin\!\left(\dfrac{5\pi}{6}\right)$ and
    $\cos\!\left(\dfrac{5\pi}{6}\right)$.
    \part[2] An arc on a circle of radius 10 cm subtends a central angle of
    $\dfrac{3\pi}{4}$. Find the arc length.
    \part[2] Find all values of $x \in [0, 2\pi]$ where $\sin x = -\dfrac{1}{2}$.
\end{parts}
\begin{solution}[9cm]
\textbf{(a)} $\ang{225} \times \dfrac{\pi}{180} = \dfrac{5\pi}{4}$ rad;\quad
$\dfrac{7\pi}{4} \times \dfrac{180}{\pi} = \ang{315}$

\textbf{(b)} $\dfrac{5\pi}{6} = \ang{150}$. Reference angle $\ang{30}$, QII.
$\sin\!\left(\dfrac{5\pi}{6}\right) = \dfrac{1}{2}$;\quad
$\cos\!\left(\dfrac{5\pi}{6}\right) = -\dfrac{\sqrt{3}}{2}$

\textbf{(c)} $a = r\theta = 10 \times \dfrac{3\pi}{4} = \dfrac{30\pi}{4} = \dfrac{15\pi}{2}$ cm

\textbf{(d)} Reference angle: $\sin^{-1}\!\left(\dfrac{1}{2}\right) = \dfrac{\pi}{6}$.
$\sin x < 0$ in QIII and QIV:
$x = \pi + \dfrac{\pi}{6} = \dfrac{7\pi}{6}$ and $x = 2\pi - \dfrac{\pi}{6} = \dfrac{11\pi}{6}$
\end{solution}

\question For each function, state the amplitude, period, phase shift, and vertical
shift. Then sketch one full cycle.
\begin{parts}
    \part[4] $y = 3\sin\!\left(2x - \dfrac{\pi}{2}\right) + 1$
    \part[3] $y = -2\cos\!\left(\dfrac{x}{2} + \pi\right)$
\end{parts}
\begin{solution}[9cm]
\textbf{(a)} Rewrite: $y = 3\sin\!\left(2\!\left(x - \dfrac{\pi}{4}\right)\right) + 1$

Amplitude: $|a| = \mathbf{3}$;\quad Period: $\dfrac{2\pi}{k} = \dfrac{2\pi}{2} = \mathbf{\pi}$

Phase shift: $\mathbf{\dfrac{\pi}{4}}$ to the right;\quad Vertical shift: \textbf{1 up}

Range: $[-2, 4]$.
Cycle starts at $x = \dfrac{\pi}{4}$, max at $x = \dfrac{\pi}{4} + \dfrac{\pi}{4} = \dfrac{\pi}{2}$,
ends at $x = \dfrac{\pi}{4} + \pi = \dfrac{5\pi}{4}$.

\textbf{(b)} Rewrite: $y = -2\cos\!\left(\dfrac{1}{2}(x + 2\pi)\right)$

Amplitude: $\mathbf{2}$;\quad Period: $\dfrac{2\pi}{1/2} = \mathbf{4\pi}$

Phase shift: $\mathbf{2\pi}$ to the left;\quad Vertical shift: \textbf{none} (0)

Reflected in $x$-axis. Range: $[-2, 2]$.
\end{solution}

\question Determine the equation of a sinusoidal function given each description.
\begin{parts}
    \part[3] A sine function with amplitude 5, period $4\pi$, phase shift $\dfrac{\pi}{3}$
    to the right, and vertical shift 2 down.
    \part[3] A cosine function whose maximum value is 7 and minimum value is $-1$,
    with period $\pi$ and no phase shift.
\end{parts}
\begin{solution}[8cm]
\textbf{(a)} $a = 5$;\; $k = \dfrac{2\pi}{4\pi} = \dfrac{1}{2}$;\; $d = \dfrac{\pi}{3}$;\; $c = -2$
\[
y = 5\sin\!\left(\tfrac{1}{2}\!\left(x - \tfrac{\pi}{3}\right)\right) - 2
\]

\textbf{(b)} Amplitude $= \dfrac{7-(-1)}{2} = 4$;\quad Vertical shift $= \dfrac{7+(-1)}{2} = 3$

$k = \dfrac{2\pi}{\pi} = 2$;\quad $d = 0$
\[
y = 4\cos(2x) + 3
\]
\end{solution}

\question A Ferris wheel has a diameter of 20 m. The centre is 12 m above the ground.
The wheel completes one full revolution every 40 seconds. A rider starts at the
lowest point.
\begin{parts}
    \part[2] State the amplitude, period, and vertical shift of the height function.
    \part[3] Write a sinusoidal equation $h(t)$ for the rider's height above the
    ground as a function of time $t$ in seconds. (Use a cosine function, noting
    the rider starts at the minimum.)
    \part[2] Find the rider's height at $t = 10$ s and $t = 25$ s.
\end{parts}
\begin{solution}[9cm]
\textbf{(a)} Amplitude $= 10$ m (radius); Period $= 40$ s;
Vertical shift $= 12$ m (centre height); min height $= 2$ m, max $= 22$ m.

\textbf{(b)} Rider starts at minimum $\Rightarrow$ use $-\cos$:
\[
h(t) = -10\cos\!\left(\frac{2\pi}{40}t\right) + 12 = -10\cos\!\left(\frac{\pi}{20}t\right) + 12
\]

\textbf{(c)} $h(10) = -10\cos\!\left(\dfrac{\pi}{2}\right) + 12 = -10(0) + 12 = \mathbf{12 \text{ m}}$
(at centre height, moving up)

$h(25) = -10\cos\!\left(\dfrac{25\pi}{20}\right) + 12 = -10\cos\!\left(\dfrac{5\pi}{4}\right) + 12$
$= -10\!\left(-\dfrac{\sqrt{2}}{2}\right) + 12 = 5\sqrt{2} + 12 \approx \mathbf{19.1 \text{ m}}$
\end{solution}

\end{questions}
\end{document}
